The degenerating direction is null for the limit form (Proposition \claimref{singularity.degeneration-exactly}(3)). This is the type-N datum of \cite[Def.~3.1, Prop.~3.3]{ButlerT1} in a degenerate setting that the root paper's taxonomy does not itself cover: there, type N is a null vector of a nondegenerate form, which acquires its annihilating idempotent through a hyperbolic partner, whereas here $\er$ lies in the radical of $\hat g$ and has no partner at all. The root paper is explicit that its algorithm assumes nondegeneracy and that the degenerate case is the sheaf paper's \cite[Rem.~4.15, Rem.~8.4]{ButlerT1}; the warrant for treating the radical direction as type N and applying Rule 1 fibrewise is therefore the localisation theorem of \cite[Thm.~4.4, Cor.~4.5]{ButlerT4} in the Schwarzschild form \cite[Thms 8.7, 8.12]{ButlerT4}, which Step 3 invokes, and which supplies the hyperbolic partner by completion exactly on $E$. The absence of type C is the classification: the mod-8 arithmetic and the classification $\Cl(1,3) \cong \Mat_2(\HH)$ are classical \cite{Porteous1995}; the existence criterion for type C ($p - q \equiv 1 \bmod 4$) is \cite[Thm.~3.6]{ButlerT1}, and $6 \not\equiv 1 \bmod 4$; the root paper's own worked signature for $\Cl(1,3)$ is \cite[Ex.~3.12]{ButlerT1}. That type P imposes nothing is \cite[Prop.~6.1]{ButlerT1}: the identity assignment repairs every idempotent at cost zero. The tower statement is the frozen clock of T3 \cite{ButlerT3}. The script checks the mod-8 datum and the centre; the type-N diagnosis itself is a citation, not a computation.
